\begin{ouverture}[algorithmique]{5}{Égalités, inconnues et régularités}{Partie 1 · Nombres entiers et calculs}
{Signe égal · Nombres cachés · Suites · Programmes de calcul}
\objectif{Compléter des égalités à trous en raisonnant.}
\objectif{Désigner un nombre inconnu par un symbole ou une lettre.}
\objectif{Poursuivre des suites et exécuter des programmes de calcul.}
\prerequis{Les quatre opérations sur les entiers.}
\end{ouverture}
\pagestyle{otMain}\essentiel\fontsize{11}{15}\selectfont
\phantomsection\label{cours-5-1}
\essSection{1}{Des égalités et des nombres cachés}
\begin{definition}Le signe $=$ relie deux écritures qui désignent le même nombre. L’égalité $8+6=9+5$ est vraie : chaque côté vaut 14. Le signe $=$ ne signifie pas seulement « voici le résultat ».\end{definition}
\begin{savoirfaire}[Compléter une égalité à trous]Pour compléter $178-\square=6\times8$, calculer d’abord le côté connu : $6\times8=48$. Puis chercher le nombre retiré à 178 pour obtenir 48 : $178-48=130$. Donc $\square=130$.\end{savoirfaire}
\begin{exemple}Un symbole peut cacher un nombre à trouver : si $\triangle\times4+3=31$, alors $\triangle\times4=28$, donc $\triangle=7$. On remonte le calcul à l’envers, opération par opération.\end{exemple}
\begin{attention}Un même symbole cache le même nombre dans toute l’égalité. Dans $\triangle+\triangle+5=19$, les deux triangles valent ensemble 14, donc chacun vaut 7.\end{attention}
\clearpage
\phantomsection\label{cours-5-2}
\essSection{2}{Suites et programmes de calcul}
\begin{definition}Une suite de nombres suit une règle. Dans $3 ;\ 10 ;\ 17 ;\ 24$, chaque terme vaut le précédent plus 7. Une règle peut combiner deux actions : dans $80 ;\ 85 ;\ 83 ;\ 88 ;\ 86$, on ajoute 5 puis on retire 2, en alternance.\end{definition}
\label{fig-cours-5-2-prog-cours}\figureGeo{prog-cours}

\begin{exemple}Exécuter le programme avec 6 : $6+2=8$, puis $8\times4=32$. Le programme affiche 32. Avec 10 : $12\times4=48$. Chaque nombre choisi donne sa propre sortie.\end{exemple}
\begin{savoirfaire}[Lire une suite de motifs]Dans une suite de motifs, compter ce qui change d’une étape à la suivante : si chaque étape ajoute 3 carreaux, l’étape suivante se déduit sans tout redessiner. Noter les nombres d’une étape à l’autre pour voir la régularité.\end{savoirfaire}

\clearpage\rubrique{Découverte}
\begin{blocLivre}{ex-5-1}\begin{exercice}[1]{Vrai ou faux ?}
L’égalité $8+6=9+5$ est-elle vraie ? Et $7\times4=30-3$ ?
\end{exercice}\end{blocLivre}
\begin{blocLivre}{ex-5-2}\begin{exercice}[1]{Trous additifs}
Compléter : $47+\square=92$ et $\square-28=35$.
\end{exercice}\end{blocLivre}
\begin{blocLivre}{ex-5-3}\begin{exercice}[1]{Trous multiplicatifs}
Compléter : $7\times\square=84$, puis $178-\square=6\times8$.
\end{exercice}\end{blocLivre}
\begin{blocLivre}{ex-5-4}\begin{exercice}[1]{La machine à nombres}
Exécuter le programme de calcul avec 6, puis avec 10 : choisir un nombre ; ajouter 2 ; multiplier par 4.
\label{fig-ex-5-4}\figureGeo{prog-4}
\end{exercice}\end{blocLivre}
\begin{blocLivre}{ex-5-5}\begin{exercice}[1]{Poursuivre une suite}
Donner les deux termes suivants de la suite $3 ;\ 10 ;\ 17 ;\ 24 ;\ \ldots$ et énoncer la règle.
\end{exercice}\end{blocLivre}

\rubrique{Entraînement}
\begin{blocLivre}{ex-5-6}\begin{exercice}[2]{Deux règles différentes}
Trouver la règle de chaque suite et donner le terme suivant : $5 ;\ 10 ;\ 20 ;\ 40 ;\ \ldots$ puis $5 ;\ 10 ;\ 15 ;\ 20 ;\ \ldots$
\end{exercice}\end{blocLivre}
\begin{blocLivre}{ex-5-7}\begin{exercice}[2]{Remonter la machine}
Un programme multiplie par 3 puis retire 4. Le résultat affiché est 20. Quel nombre a été choisi ?
\end{exercice}\end{blocLivre}
\begin{blocLivre}{ex-5-8}\begin{exercice}[2]{La suite qui zigzague}
Donner les deux termes suivants : $80 ;\ 85 ;\ 83 ;\ 88 ;\ 86 ;\ 91 ;\ 89 ;\ \ldots$
\end{exercice}\end{blocLivre}
\begin{blocLivre}{ex-5-9}\begin{exercice}[2]{Deux symboles}
Trouver les nombres cachés : $\triangle+\triangle+5=19$, puis $\triangle+\bullet=12$.
\end{exercice}\end{blocLivre}
\begin{blocLivre}{ex-5-10}\begin{exercice}[2]{Le nombre caché}
Trouver le nombre caché : $\square\times4+3=31$.
\end{exercice}\end{blocLivre}

\rubrique{Problèmes}
\begin{blocLivre}{ex-5-11}\begin{exercice}[3]{Les tee-shirts}
Un tee-shirt coûte 12 euros et la livraison coûte 5 euros, quel que soit le nombre acheté.\begin{questions}
\item Calculer le prix de 4 tee-shirts livrés.
\item Écrire le prix à payer pour $N$ tee-shirts livrés.
\end{questions}
\end{exercice}\end{blocLivre}
\begin{blocLivre}{ex-5-12}\begin{exercice}[3]{Ciseaux et stylos}
Deux paires de ciseaux identiques et trois stylos à 2 euros coûtent 20 euros en tout. Quel est le prix d’une paire de ciseaux ?
\end{exercice}\end{blocLivre}
\begin{blocLivre}{ex-5-13}\begin{exercice}[3]{Des carrés d’allumettes}
Une rangée de carrés est construite avec des allumettes : le premier carré en demande 4, chaque carré suivant en demande 3 de plus. Combien d’allumettes pour une rangée de 5 carrés ? Et pour 10 carrés ?
\label{fig-ex-5-13}\figureGeo{motifs-13}
\end{exercice}\end{blocLivre}
\begin{blocLivre}{ex-5-14}\begin{exercice}[3]{Le motif en croix}
Une croix de carreaux grandit à chaque étape : 5 carreaux à l’étape 1, puis 4 carreaux de plus à chaque étape. Combien de carreaux à l’étape 4 ? À quelle étape utilise-t-on exactement 29 carreaux ?
\label{fig-ex-5-14}\figureGeo{motifs-14}
\end{exercice}\end{blocLivre}
\begin{blocLivre}{ex-5-15}\begin{exercice}[3]{Inventer un programme}
Proposer un programme de calcul en deux étapes qui transforme 5 en 26. Les programmes « multiplier par 5 puis ajouter 1 » et « ajouter 1 puis multiplier par 5 » donnent-ils le même résultat ?
\end{exercice}\end{blocLivre}

\clearpage
